\documentclass{article}
\usepackage[T1]{fontenc}
\usepackage{lmodern}
\usepackage{graphicx}
\usepackage{amsmath,amssymb}
\usepackage[a4paper,margin=2cm]{geometry}
\usepackage[absolute,overlay]{textpos}
\setlength{\TPHorizModule}{1mm}
\setlength{\TPVertModule}{1mm}
\usepackage[indonesian]{babel}
\usepackage{hyperref}
\setlength{\emergencystretch}{2em}
% unit: Halaman 54

\begin{document}

% === Halaman 54 (llm) ===

\section*{Galois Field $\text{GF}(2^m)$}

\begin{itemize}
    \item Disebut juga medan berhingga biner.
    \item $\text{GF}(2^m)$ atau $\text{F}_{2^m}$ adalah ruang vektor berdimensi $m$ pada $\text{GF}(2)$. Setiap elemen di dalam $\text{GF}(2^m)$ adalah integer dalam representasi biner sepanjang maksimal $m$ bit.\
    Contoh: $\text{GF}(2^4)$, $\text{GF}(2^8)$, dst
    \item String biner $b_{m-1}b_{m-2} \dots b_1b_0$, $b_i \in \{0,1\}$, dapat dinyatakan sebagai polinom
    \[ b_{m-1}x^{m-1} + b_{m-2}x^{m-2} + \dots + b_1x + b_0 \]
    \item Jadi, setiap $a \in \text{GF}(2^m)$ dapat dinyatakan sebagai
    \[ b_{m-1}x^{m-1} + b_{m-2}x^{m-2} + \dots + b_1x + b_0 \]
    \item Contoh: 1101 di dalam $\text{GF}(2^4)$ dapat dinyatakan sebagai polinom $x^3 + x^2 + 1$\\
    01010111 di dalam $\text{GF}(2^8)$ dapat dinyatakan sebagai polinom $x^6 + x^4 + x^2 + x + 1$
\end{itemize}

\end{document}
